\documentclass[aps,pra,reprint,superscriptaddress,nofootinbib,longbibliography]{revtex4-2}

\usepackage{amsmath,amssymb,bm,mathtools}
\usepackage{graphicx}

\newcommand{\id}{\mathbb I_2}
\newcommand{\R}{\mathbb R}
\newcommand{\bmnu}{\bm{\mu}}
\newcommand{\bn}{\bm n}
\newcommand{\bu}{\bm u}
\newcommand{\bx}{\bm x}

\begin{document}

\title{Exact joint-measurability threshold for four coplanar qubit measurements}

\author{Xiaojian Zeng}
\email{xiaojianzeng@foxmail.com}
\affiliation{PBC School of Finance, Tsinghua University, 100190}

\date{September 12, 2026}

\begin{abstract}
Four projective spin measurements on a qubit become jointly measurable after
sufficient common depolarization.  For four distinct coplanar axes, a known
construction gives a visibility threshold in terms of the successive angular
gaps, but its optimality for unequal gaps has previously been supported only
numerically.  We prove that the construction is optimal throughout the strict
ordered semicircle domain.  The proof evaluates the unrestricted
four-measurement Fourier criterion exactly.  The vector part of the known
parent measurement supplies a feasible four-support primal point, while an
explicit affine dual function has matching value.  Its only nontrivial
feasibility conditions reduce to an elementary inequality for three ordered
planar unit vectors.  Consequently, if
$0=\alpha_1<\alpha_2<\alpha_3<\alpha_4<\pi$, the maximum common visibility is
\(
 [\sum_{j=2}^{4}\sin((\alpha_j-\alpha_{j-1})/2)
 +\cos(\alpha_4/2)]^{-1}
\).
The result assumes unbiased binary qubit effects and an unrestricted parent
POVM; it makes no claim for unequal visibilities, biased effects, or
noncoplanar measurements.
\end{abstract}

\maketitle

\section{Introduction}

Joint measurability asks whether the statistics of several measurements can
be recovered as marginals of a single parent positive-operator-valued measure
(POVM).  For noisy qubit measurements this question gives a precise
operational measure of how much common white noise is needed before one
simultaneous readout can simulate all measurement choices
\cite{UolaEtAl2016,AndrejicKunjwal2020,GrinkoUola2025}.

For coplanar unbiased binary qubit measurements of a common visibility,
Andrejic and Kunjwal constructed a parent POVM whose admissible visibility is
set by a polygonal norm sum \cite{AndrejicKunjwal2020}.  They conjectured that
this sufficient condition is also necessary.  Grinko and Uola subsequently
gave a necessary-and-sufficient Fourier characterization for any finite
family of unbiased binary qubit measurements \cite{GrinkoUola2025}; however,
using it still requires a nontrivial norm minimization.  Most recently, Porto
\emph{et al.} found strong numerical evidence that the polygonal bound is
tight for planar projective measurements \cite{PortoEtAl2026}.

Here we settle the four-measurement specialization analytically.  We evaluate
all eight norm channels and retain all four free triple-index Fourier vectors.
The proof is a matching primal--dual certificate and therefore does not
restrict a putative parent POVM to the support pattern of the known
construction.  The result fixes the exact common depolarizing visibility for
four arbitrary distinct coplanar qubit axes.  It does not address the
arbitrary-number conjecture.

\section{Model and exact statement}
\label{sec:model}

Let $\bm\sigma=(\sigma_x,\sigma_y,\sigma_z)$, with
$\operatorname{tr}(\sigma_k\sigma_l)=2\delta_{kl}$.  Common rotations,
independent outcome relabelings, and permutation of four distinct coplanar
projective axes allow representatives
\begin{equation}
  \bn_j=(\cos\alpha_j,\sin\alpha_j,0),\qquad
  0=\alpha_1<\alpha_2<\alpha_3<\alpha_4<\pi .
  \label{eq:ordered-domain}
\end{equation}
For outcome $a\in\{-1,+1\}$ and common visibility $\eta\in[0,1]$, define
\begin{equation}
 E_{a|j}^{(\eta)}
 =\frac12\bigl(\id+a\eta\,\bn_j\mathbin{\cdot}\bm\sigma\bigr).
 \label{eq:noisy-effects}
\end{equation}
Thus $1-\eta$ is the depolarizing-noise fraction.  The four POVMs are jointly
measurable when there are sixteen positive semidefinite effects
$G_{\bm s}$, $\bm s\in\{-1,+1\}^4$, such that
\begin{align}
 \sum_{\bm s}G_{\bm s}&=\id,
 \nonumber\\
 \sum_{\bm s:s_j=a}G_{\bm s}&=E_{a|j}^{(\eta)}
 \quad (j=1,\ldots,4;\ a=\pm1).
 \label{eq:parent}
\end{align}
No rank, planarity, or support restriction is imposed on the parent; zero
effects are allowed.

Introduce
\begin{equation}
 D=\sum_{j=2}^{4}\sin\frac{\alpha_j-\alpha_{j-1}}2
   +\cos\frac{\alpha_4}{2}.
 \label{eq:D-trig}
\end{equation}
The main result is the following.

\textbf{Theorem.} \emph{Under Eq.~\eqref{eq:ordered-domain}, the effects in
Eq.~\eqref{eq:noisy-effects} admit an unrestricted parent POVM if and only if
\begin{equation}
 \eta D\leq 1.
 \label{eq:threshold-criterion}
\end{equation}
Consequently their maximum common visibility is
\begin{equation}
 \eta_*(\alpha_2,\alpha_3,\alpha_4)=D^{-1}.
 \label{eq:eta-star}
\end{equation}}

Figure~\ref{fig:geometry} shows the two pieces of planar geometry used below.
The centrally symmetric polygon is useful for interpreting $D$, whereas the
ordered arc of four auxiliary unit vectors controls dual feasibility.

\begin{figure*}[t]
 \centering
 \includegraphics[width=0.84\textwidth]{manuscript/figures/geometry.pdf}
 \caption{Geometry of the exact certificate for a representative strict
 ordered configuration.  (a) The highlighted path from $\bn_1$ through
 $\bn_4$ to $-\bn_1$ has length $2D$; central symmetry therefore gives
 perimeter $4D$ for $\operatorname{conv}\{\pm\bn_j\}$.  (b) The four unit
 directions $\bu_1,\ldots,\bu_4$ lie in angular order on an arc shorter than
 $\pi$.  The displayed endpoints-minus-middle vector illustrates the lemma
 that bounds every off-support dual value by the unit disk.  The figure is a
 schematic, not numerical evidence.  It was produced with Matplotlib code
 developed with assistance from OpenAI Codex and checked against the analytic
 definitions.}
 \label{fig:geometry}
\end{figure*}

\section{Fourier compatibility criterion}
\label{sec:fourier}

We use the necessary-and-sufficient criterion of Grinko and Uola
\cite{GrinkoUola2025}.  On the half cube
\begin{equation}
 \mathcal D_4=\{(1,\mu_2,\mu_3,\mu_4):\mu_k\in\{-1,+1\}\},
\end{equation}
write
\begin{equation}
 \bm X=(\bx_{123},\bx_{124},\bx_{134},\bx_{234})\in(\R^3)^4
\end{equation}
and define
\begin{align}
 q(\bn_1,\ldots,\bn_4)
  &=\min_{\bm X\in(\R^3)^4}
    \sum_{\bmnu\in\mathcal D_4}\|\bm r(\bmnu)\|,
 \label{eq:q-def}\\[-1ex]
 \bm r(\bmnu)={}&\bn_1+\mu_2\bn_2+\mu_3\bn_3+\mu_4\bn_4
 \nonumber\\
 &+\mu_2\mu_3\bx_{123}+\mu_2\mu_4\bx_{124}\nonumber\\
 &+\mu_3\mu_4\bx_{134}
 +\mu_2\mu_3\mu_4\bx_{234}.
 \label{eq:r-def}
\end{align}
All four vectors in the minimization are unrestricted elements of $\R^3$.
For the common-visibility effects, joint measurability is equivalent to the
same minimum with fixed vectors $\eta\bn_j$ being at most $8$.  Positive
homogeneity of Eq.~\eqref{eq:q-def} gives
\begin{equation}
 q(\eta\bn_1,\ldots,\eta\bn_4)
 =\eta q(\bn_1,\ldots,\bn_4)
 \label{eq:homogeneity}
\end{equation}
for $\eta\geq0$.  It therefore remains to prove $q=8D$ at unit visibility.

\section{Matching exact certificates}
\label{sec:certificate}

Define the positive half-chord lengths and corresponding unit directions
\begin{align}
 d_j&=\frac{\|\bn_j-\bn_{j+1}\|}{2},&
 \bu_j&=\frac{\bn_j-\bn_{j+1}}{2d_j}
 &&(j=1,2,3),
 \label{eq:chord-defs}\\
 c&=\frac{\|\bn_1+\bn_4\|}{2},&
 \bu_4&=\frac{\bn_1+\bn_4}{2c}.
 \label{eq:endpoint-defs}
\end{align}
Then $D=d_1+d_2+d_3+c$.

\subsection{Primal certificate}

Define a vector-valued function $R$ on $\mathcal D_4$ by the four nonzero
values
\begin{table}[b]
 \caption{Support of the primal function $R$.}
 \label{tab:primal}
 \begin{ruledtabular}
 \begin{tabular}{cc}
 $(\mu_2,\mu_3,\mu_4)$ & $R(\bmnu)$ \\
 \hline
 $(-,-,-)$ & $8d_1\bu_1=4(\bn_1-\bn_2)$ \\
 $(+,-,-)$ & $8d_2\bu_2=4(\bn_2-\bn_3)$ \\
 $(+,+,-)$ & $8d_3\bu_3=4(\bn_3-\bn_4)$ \\
 $(+,+,+)$ & $8c\bu_4=4(\bn_1+\bn_4)$
 \end{tabular}
 \end{ruledtabular}
\end{table}
and set $R=0$ at the other four vertices.  To see directly that this is a
feasible point of Eq.~\eqref{eq:q-def}, denote the four displayed values by
$A,B,C,F$ in their table order.  Their normalized Walsh sums are
\begin{align}
 \tfrac18(A+B+C+F)&=\bn_1,\nonumber\\
 \tfrac18(-A+B+C+F)&=\bn_2,\nonumber\\
 \tfrac18(-A-B+C+F)&=\bn_3,\nonumber\\
 \tfrac18(-A-B-C+F)&=\bn_4.
 \label{eq:low-walsh}
\end{align}
These are exactly the coefficients of the constant and the three singleton
characters $1,\mu_2,\mu_3,\mu_4$.  The remaining Walsh characters are
$\mu_2\mu_3$, $\mu_2\mu_4$, $\mu_3\mu_4$, and
$\mu_2\mu_3\mu_4$; their four coefficients can therefore be taken as the
four free vectors in Eq.~\eqref{eq:r-def}.  Thus $\bm r=R$ is feasible and
\begin{equation}
 q\leq\sum_{\bmnu}\|R(\bmnu)\|
 =8(d_1+d_2+d_3+c)=8D.
 \label{eq:primal-bound}
\end{equation}
This feasible point is the Walsh transform of the vector part of the known
parent construction \cite{AndrejicKunjwal2020}; it is not a support
restriction on the minimization.

\subsection{Planar unit-vector lemma}

Let $\bu(\theta)=(\cos\theta,\sin\theta,0)$.  If
$a\leq b\leq c_0$ and $c_0-a\leq\pi$, set $p=b-a$ and $s=c_0-b$.
An elementary expansion gives
\begin{align}
 1-\|\bu(a)+\bu(c_0)-\bu(b)\|^2
 &=8\cos\frac{p+s}{2}\sin\frac p2\sin\frac s2
 \geq0.
 \label{eq:unit-lemma}
\end{align}
Hence every endpoints-minus-middle combination of three unit vectors ordered
on a closed semicircle has norm at most one.

\subsection{Dual certificate}

For every half-cube vertex define
\begin{align}
 \bm Y(\bmnu)={}&-\frac{\bu_1+\bu_4}{2}
 +\frac{\mu_2}{2}(\bu_1-\bu_2)\nonumber\\
 &+\frac{\mu_3}{2}(\bu_2-\bu_3)
 +\frac{\mu_4}{2}(\bu_3-\bu_4).
 \label{eq:dual}
\end{align}
Because $\bm Y$ contains only the constant and singleton characters,
Boolean-character orthogonality gives
\begin{equation}
 \sum_{\bmnu\in\mathcal D_4}\chi_S(\bmnu)\bm Y(\bmnu)=0,
 \quad S\in\{123,124,134,234\},
 \label{eq:dual-orthogonality}
\end{equation}
where the four $\chi_S$ are the optimization characters appearing in
Eq.~\eqref{eq:r-def}.

The eight values of $\bm Y$ are listed in Table~\ref{tab:dual}.  At the four
support vertices of $R$ they are $-\bu_1,-\bu_2,-\bu_3,-\bu_4$,
respectively.  To control the other four, note that the polar angles of the
auxiliary directions are
\begin{align}
 \beta_1&=\frac{\alpha_2}{2}-\frac\pi2,&
 \beta_2&=\frac{\alpha_2+\alpha_3}{2}-\frac\pi2,\nonumber\\
 \beta_3&=\frac{\alpha_3+\alpha_4}{2}-\frac\pi2,&
 \beta_4&=\frac{\alpha_4}{2}.
 \label{eq:beta}
\end{align}
They satisfy $\beta_1<\beta_2<\beta_3<\beta_4$ and
\begin{equation}
 \beta_4-\beta_1=\frac{\pi+\alpha_4-\alpha_2}{2}<\pi.
 \label{eq:beta-span}
\end{equation}
Every off-support entry in Table~\ref{tab:dual} is therefore covered by
Eq.~\eqref{eq:unit-lemma}.  The support entries have norm one, so
\begin{equation}
 \|\bm Y(\bmnu)\|\leq1\qquad(\bmnu\in\mathcal D_4).
 \label{eq:dual-feasible}
\end{equation}

\begin{table}[t]
 \caption{The dual function.  The first, fifth, seventh, and eighth rows are
 the support vertices of $R$ in Table~\ref{tab:primal}.}
 \label{tab:dual}
 \begin{ruledtabular}
 \begin{tabular}{cc}
 $(\mu_2,\mu_3,\mu_4)$ & $\bm Y(\bmnu)$ \\
 \hline
 $(-,-,-)$ & $-\bu_1$ \\
 $(-,-,+)$ & $-(\bu_1+\bu_4-\bu_3)$ \\
 $(-,+,-)$ & $-(\bu_1+\bu_3-\bu_2)$ \\
 $(-,+,+)$ & $-(\bu_1+\bu_4-\bu_2)$ \\
 $(+,-,-)$ & $-\bu_2$ \\
 $(+,-,+)$ & $-(\bu_2+\bu_4-\bu_3)$ \\
 $(+,+,-)$ & $-\bu_3$ \\
 $(+,+,+)$ & $-\bu_4$
 \end{tabular}
 \end{ruledtabular}
\end{table}

For arbitrary values of all four free vectors in Eq.~\eqref{eq:r-def},
Cauchy--Schwarz, Eqs.~\eqref{eq:dual-orthogonality} and
\eqref{eq:dual-feasible} imply
\begin{align}
 \sum_{\bmnu}\|\bm r(\bmnu)\|
 &\geq-\sum_{\bmnu}\bm Y(\bmnu)\mathbin{\cdot}\bm r(\bmnu)\nonumber\\
 &=-\sum_{\bmnu}\bm Y(\bmnu)\mathbin{\cdot}
 (\bn_1+\mu_2\bn_2+\mu_3\bn_3+\mu_4\bn_4).
 \label{eq:weak-duality}
\end{align}
The fixed low-order coefficients in Eq.~\eqref{eq:low-walsh} allow the last
line to be evaluated with $R$.  At each nonzero vertex,
$\bm Y=-\bu_j$ and $R$ is a positive multiple of $\bu_j$, giving
\begin{equation}
 -\sum_{\bmnu}\bm Y(\bmnu)\mathbin{\cdot}R(\bmnu)
 =8d_1+8d_2+8d_3+8c=8D.
 \label{eq:dual-value}
\end{equation}
Equations~\eqref{eq:primal-bound} and \eqref{eq:dual-value} prove the global,
unrestricted equality
\begin{equation}
 \boxed{q(\bn_1,\bn_2,\bn_3,\bn_4)=8D}.
 \label{eq:q-result}
\end{equation}
This conclusion uses only weak duality; no numerical strong-duality
assumption is involved.  Combining Eqs.~\eqref{eq:homogeneity} and
\eqref{eq:q-result} with the Fourier compatibility criterion proves
Eq.~\eqref{eq:threshold-criterion}.

\section{Physical range and checks}
\label{sec:checks}

The threshold lies in the physical visibility interval.  Indeed, $2D$ is
the total length of the four vectors
\begin{equation}
 \bn_1-\bn_2,\quad \bn_2-\bn_3,\quad
 \bn_3-\bn_4,\quad \bn_4+\bn_1,
\end{equation}
whose sum is $2\bn_1$.  The triangle inequality gives $D\geq1$, and equality
would require collinearity incompatible with the strict ordering in
Eq.~\eqref{eq:ordered-domain}.  Therefore $D>1$ and $0<\eta_*<1$.

For the symmetric angles
$(0,\pi/4,\pi/2,3\pi/4)$, Eq.~\eqref{eq:D-trig} gives
$D=4\sin(\pi/8)$ and recovers the known exact threshold
$[4\sin(\pi/8)]^{-1}$ \cite{UolaEtAl2016,AndrejicKunjwal2020}.
The unequal-gap proof is exact and does not use a limiting expansion near
this configuration.

Computations were used only to attack the certificate, not to establish the
theorem.  The accompanying verification program performs exact symbolic
Walsh and trigonometric checks, solves the full four-vector eight-channel
second-order-cone program, and separately optimizes sixteen unrestricted
$2\times2$ Hermitian parent effects.  Across 128 seeded hold-out angle sets,
the maximum observed $|q_{\mathrm{SOCP}}-8D|$ was
$2.68\times10^{-9}$ and the largest reported solver primal--dual gap was
$5.72\times10^{-10}$.  Across eight named direct-parent checks, the computed
visibility agreed with $D^{-1}$ within $7.9\times10^{-9}$.  Residuals,
tolerances, environment versions, and all sampled cases are recorded in the
reproducibility package.

\subsection{AI-assisted research and verification}

OpenAI Codex, using GPT-6 Astra (Max reasoning) and GPT-5.6 Sol (Ultra
reasoning) model configurations, was used substantively for literature
organization, derivation discovery, code
generation and debugging, symbolic and numerical checks, adversarial review,
and manuscript drafting.  The generated results were cross-checked within a
staged workflow by independent source reconstruction, exact Walsh
enumeration, symbolic factorization, and separately implemented numerical
optimizations.  These automated checks do not replace author judgment or
accountability.  This local draft requires the named author's personal
scientific review, correction where needed, confirmation of the disclosure,
and approval before any submission.


\section{Scope and conclusion}

The exact formula applies to one qubit, four distinct coplanar projective
axes represented in one open semicircle, binary unbiased effects, and one
common depolarizing visibility.  The parent POVM is unrestricted in
$\mathbb C^2$ and may contain zero effects.  Coincident or antipodal boundary
points are outside the stated strict domain, although the formula has
well-defined one-sided limits.  Unequal visibilities, biased effects,
noncoplanar directions, more or fewer than four measurements,
higher-dimensional or nonbinary measurements, sequential nondisturbance,
and multicopy notions require separate analyses.

Within this domain, the affine dual certificate supplies the missing
necessity direction of the known parent construction.  Thus the polygonal
noise bound is not merely sufficient: for four axes it is the exact common
white-noise threshold for simulation by a single parent measurement.

\section*{Data availability}

No experimental data were used.  The source package accompanying this draft
contains the verification code, frozen JSON output, plain-text log, pinned
environment versions, and regeneration instructions.  The analytic proof is
self-contained and does not depend on the numerical calculations.

\begin{acknowledgments}
OpenAI Codex also assisted substantively with manuscript organization and
drafting.  The research assistance and automated checks are described above;
the named author's personal review, correction, and approval remain pending.
\end{acknowledgments}

\bibliography{manuscript/references}

\end{document}
